% Using InitechBase, Appendix A: Keys at the Dot Prompt (annex Sec 8.2
% UA, "Cursor Control Keys", thinned to the keys the shipped dot prompt
% reads).  Sources: os/samir/pal/pal_milton.c milton_terminal_line (with
% the status bar: Enter ends the line, Backspace erases, printable
% characters 32-126 are taken, other keys ignored; 254 characters) and
% samir_line_read over the InitechDOS console (without the status bar);
% cmd/proc.c (WAIT: one key, Enter = ""); cmd/mutate.c m_confirm_file (Y).
\chapter{Keys at the Dot Prompt}
\Gates{test-samir-input, test-samir-status}
\NoDictionaryHeads
InitechBase reads each command, and each answer to \cmd{ACCEPT} and
\cmd{INPUT}, as one line typed at the keyboard.

\noindent\hspace*{-\imHang}\begin{tabular}{@{}p{34mm}p{100mm}@{}}
\hline
\textsf{\textbf{Key}} & \textsf{\textbf{Action}}\\
\hline
\key{Enter} & Ends the line.  InitechBase carries out the command, or takes the
  answer.\\
\key{Backspace} & Erases the character before the cursor.\\
Letters, digits, punctuation, \key{Space} & Typed into the line, up to 254
  characters.\\
\hline
\end{tabular}

While the status bar is displayed (SET STATUS ON), InitechBase reads the line
itself and takes only the keys above; any other key is ignored.  While it is
not, the line is read through InitechDOS, and the line-editing keys of the
InitechDOS prompt apply (\doctitle{InitechDOS User's Guide}, ``The InitechDOS
Editing Keys'').

\section{Answering WAIT}
\cmd{WAIT} takes the first key pressed, without \key{Enter}.  \key{Enter}
itself gives an empty string.

\section{Answering ZAP}
While SET SAFETY is ON, \cmd{ZAP} asks \texttt{(Y/N)}.  Type \type{Y} and press
\key{Enter} to remove the records; any other answer keeps them.
